\chapter*{Прежде чем начать}

Эта книга --- о машине, которая весит полтора килограмма, потребляет двадцать ватт, как тусклая лампочка, и при этом узнаёт лица, учит языки и ловит летящий мяч. О вашем мозге. Мы будем смотреть на него так, как инженер смотрит на незнакомую плату: что за детали, как они соединены, какие сигналы бегут по проводам и что эта схема вычисляет.

Вам не понадобятся ни формулы, ни знание биологии. Понадобится любопытство и готовность представить себе провод, лампочку, выключатель. Всё остальное мы соберём по дороге.

У книги есть вторая, полная версия --- с уравнениями, моделями, задачами и ссылками на опыты. Главы в двух версиях идут в одном и том же порядке и называются по номерам одинаково. Если какая-то глава вас зацепила, откройте её полную версию: там вы найдёте, откуда взялось каждое утверждение и как его проверить самому. На сайте книги каждую главу можно переключать между «Коротко» и «Полностью» одной кнопкой.

Одно предупреждение. Про мозг известно далеко не всё, и я буду честно говорить, где кончается измеренное и начинается догадка. Это не недостаток книги, а самое интересное в ней: многие главы кончаются вопросами, на которые пока никто не знает ответа.
